\section{Evaluation and Experimental Results}

\subsection{Test Environment}
\begin{itemize}
    \item Board mạch: ESP32-S3 với phiên bản Yolo UNO.
    \item Framework: Arduino kết hợp FreeRTOS trên môi trường PlatformIO.
    \item Cảm biến: DHT11.
    \item Hiển thị: OLED SH1106 sử dụng thư viện U8g2.
    \item Network:
    \begin{itemize}
        \item AP mode: local dashboard + WebSocket
        \item STA mode: MQTT telemetry tới CoreIOT
    \end{itemize}
\end{itemize}

\subsection{Functional Test Matrix}
\begin{table}[H]
\centering
\begin{tabular}{|p{1.6cm}|p{6.8cm}|p{2.8cm}|p{3.5cm}|}
\hline
\textbf{Task} & \textbf{Nội dung kiểm thử} & \textbf{Kết quả} & \textbf{Ghi chú} \\
\hline
1 & LED single-channel có 3 profile blink và truy cập GPIO có mutex & Pass module-level & 3 profile đã hiện thực, có thể tái sử dụng cho thermal alert \\
\hline
2 & NeoPixel phản ánh môi trường theo nhiều mức và cập nhật thread-safe & Pass runtime & Có Normal/Warning/Critical, mutex bảo vệ state \\
\hline
3 & OLED hiển thị nhiều trạng thái nhiệt/ẩm và chạy theo luồng event từ semaphore & Pass runtime & Có trạng thái no-sensor/warn/alert; event từ sensor semaphore \\
\hline
4 & AP dashboard điều khiển tối thiểu 2 thiết bị, có nút hành động rõ ràng & Pass runtime & Fan/Pump + add/toggle/delete relay + save settings \\
\hline
5 & TinyML chạy on-device và phát output ra nhiều kênh & Pass runtime & Output xuất ra OLED, AP, CoreIOT; auto fan trong NORMAL/AP \\
\hline
6 & Publish telemetry lên CoreIOT trong STA mode đúng token & Pass runtime & Có retry/reconnect và log trạng thái kết nối \\
\hline
\end{tabular}
\caption{Kết quả kiểm thử chức năng theo 6 task}
\label{tab:functional-results}
\end{table}

\subsection{TinyML Evaluation}
\subsubsection{Dataset and Training Metrics}
Theo file \texttt{ml\_artifacts/tinyml\_training\_summary.json}:
\begin{itemize}
    \item Dataset: 90 mẫu nhiệt độ/độ ẩm có nhãn (3 lớp hành động).
    \item Tách tập: 72 train, 18 validation.
    \item Train accuracy: 1.0.
    \item Validation accuracy: 1.0.
    \item Số epoch train: 150.
\end{itemize}

\subsubsection{On-device Execution}
\begin{itemize}
    \item Model int8 nhúng trực tiếp vào firmware.
    \item Interpreter TFLM khởi tạo thành công với tensor arena \SI{8}{KB}.
    \item Suy luận chạy theo chu kỳ đọc cảm biến (\SI{3}{s}), sau đó cập nhật:
    \begin{itemize}
        \item \texttt{tinyml\_result}, \texttt{tinyml\_action\_id}, \texttt{tinyml\_confidence} cho cloud/AP.
        \item Dòng \texttt{AI:...} trên OLED.
        \item Logic auto fan ở mode NORMAL/AP.
    \end{itemize}
\end{itemize}

\subsubsection{Discussion on Accuracy}
Kết quả định lượng hiện có là độ chính xác ngoại tuyến từ quy trình huấn luyện. Ở mức chương trình điều khiển, nhóm đã xác nhận được tính đúng đắn của luồng suy luận thông qua nhật ký hệ thống và hành vi của thiết bị chấp hành. Để tăng tính học thuật cho phần đánh giá độ chính xác trên phần cứng, nhóm đề xuất bước tiếp theo là bổ sung bộ kiểm thử phát lại các tập mẫu có nhãn trực tiếp trên bảng mạch và tự động tính toán ma trận nhầm lẫn tại thời điểm thực thi.

\subsection{Reliability Observations}
\begin{itemize}
    \item Hệ thống giữ được kết nối nhờ cơ chế retry WiFi/MQTT theo chu kỳ.
    \item OTA subscription được reset khi mất MQTT để tránh mất luồng cập nhật firmware.
    \item Logging có mutex giúp giảm tình trạng trộn dòng khi nhiều task chạy đồng thời.
\end{itemize}
